\pdfmapfile{+cm.map}
\pdfmapfile{+lm.map}
\pdfmapfile{+symbols.map}
\documentclass[11pt]{article}
\usepackage[a4paper,margin=25mm]{geometry}
\usepackage{amsmath,amssymb,amsthm,booktabs,longtable,array}
\usepackage[T1]{fontenc}
\usepackage{lmodern,microtype,xurl}
\usepackage[hidelinks]{hyperref}
\newtheorem{theorem}{Theorem}
\newtheorem{lemma}[theorem]{Lemma}
\newtheorem{proposition}[theorem]{Proposition}
\newcommand{\ipra}{\mathsf{iPRA}}
\newcommand{\Run}{\operatorname{Run}}
\newcommand{\enc}{\operatorname{enc}}
\newcommand{\rev}{\operatorname{rev}}
\newcommand{\Val}{\operatorname{Val}}
\newcommand{\Valid}{\operatorname{Valid}}
\title{The literal physical/source bridge for PA194AX11H and PA194QAX11H}
\author{Independent arithmetic audit}
\date{14 November 2026}
\begin{document}
\maketitle
\begin{abstract}
The literal symbolic certificates for these two machines have a sound
interpretation in intuitionistic primitive recursive arithmetic.  This note
gives the required arithmetic semantics and proves each rule of the actual
interpreter, including its four list traversals, its normalization rules,
and its variable-duration composition.  Together with the finite literal
replays, these proofs establish in $\ipra$ that each machine halts on the
blank tape exactly when its supplied register program halts from its stated
initial store.  This is a physical/source theorem.  It does not identify the
register program with a PA proof enumerator, and it does not by itself prove
an equivalence with $\operatorname{Con}(\mathsf{PA})$.
\end{abstract}

\section{Objects, finite facts, and the precise conclusion}

The objects audited are the files under
\path{machines/PA194AX11H} and \path{machines/PA194QAX11H} in the
accompanying research bundle.  The implementation is
\path{verification/check_literal.py}, with
\path{verification/reference/symbolic_runtime.py},
\path{independent_runtime.py}, and \path{source_language.py}.
The first of these selects \path{HEADER_BASE=48}; the larger default in
the retained reference module is not the value used here.

\begin{center}
\begin{tabular}{lrr}
\toprule
Finite fact & AX11H & QAX11H\\
\midrule
Working controls / literal rows & 194 / 388 & 194 / 388\\
Source sites & 3671 & 3671\\
Increment / decrement sites & 1456 / 1107 & 1456 / 1107\\
Continue / break / halt sites & 1032 / 75 / 1 & 1032 / 75 / 1\\
Contracts (two per decrement) & 4778 & 4778\\
Nonhalting endpoint checks & 4777 & 4777\\
Maximum address length & 37 & 36\\
Symbolic interpreter events & 1101836 & 1130254\\
Maximum events in one contract & 403 & 403\\
Minimum constant in a duration & 3 & 3\\
Blank startup transitions & 57922 & 57922\\
\bottomrule
\end{tabular}
\end{center}

These are finite closed facts, not induction principles or assumptions about
large register values.  The full literal replay derives every contract from
the corresponding literal transition table.  The additional executable
\path{physical_ipra_audit.py} checks the complete contract index, all
4778 affine durations, and all 4777 nonhalting endpoints against both
independently resolved source graphs.  It also rederives the four block
traversals and checks their equality with the stored data.  Its receipt
records hashes of all inputs and of the verifier modules.

The dense tables are conjugate, by the supplied injective label map, to
the native tables.  The native blank start is 74, the native source boundary
is 84, and the native halting control is 254.  The symbolic interpreter
renames the boundary to 96 and retains 254 for halt.  These are names of
existing controls; no extra physical control is introduced by this proof.

Write $M$ for either literal table and $P$ for its own
\texttt{instructions.json}.  The program's initial ordinal site is 25 and
its initial logical store is $(2,0,\ldots,0)$.  We prove
\[
 \ipra\ \vdash\
 \bigl(\exists t\; M\text{ has halted on blank tape by }t\bigr)
 \ \longleftrightarrow\
 \bigl(\exists k\; P\text{ has halted from that store by }k\bigr).
 \tag{1}
\]
The theorem refers to the exact low-level register program, not an
unidentified high-level scheduler.  Its negated version follows
intuitionistically as well.

\section{Primitive recursive arithmetization}

We use the usual language of $\ipra$ with symbols for primitive recursive
functions.  All auxiliary predicates below are characteristic functions of
primitive recursive tests.  Thus an equality, implication, or Boolean
combination involving them is quantifier free.  Definitions by primitive
recursion and quantifier-free induction suffice throughout.

\subsection{Tape configurations and the halting code}

A finite-support binary tape at a scanned cell is represented by two
natural numbers $L,R$.  The low bit of $R$ is the scanned symbol; the low bit
of $L$ is the symbol immediately to its left.  Successive higher bits give
successively more distant cells.  These numbers represent the complete
infinite tape with zero exteriors: extra exterior zeros do not change a
number.  If desired, retain a signed head coordinate $h$ using any fixed
primitive recursive integer coding.  We retain it below to express exact
head displacement.

For a nonhalting configuration $(q,h,L,R)$ and a row
$\delta(q,R\bmod 2)=(w,d,q')$, the update before canonicalizing halt is
\[
\begin{array}{ll}
d=+1:&(q',h+1,\;2L+w,\;\lfloor R/2\rfloor),\\[2pt]
d=-1:&(q',h-1,\;\lfloor L/2\rfloor,\;
       (L\bmod 2)+2w+4\lfloor R/2\rfloor).
\end{array}                                                    \tag{2}
\]
If $q'$ is the native halt control, replace this configuration by the single
arithmetic code $H$.  Put $\Delta(H)=H$.  This discards tape and head only
after physical halt.  It is an arithmetic convention, not an additional
physical transition or working state.  Its halting predicate is exactly
that of the uncollapsed literal machine, by induction on the number of
literal transitions.  In particular, the CH contract need only establish
the halt control; it need not establish a particular halted tape.

Let $\Run(0,c)=c$ and
$\Run(t+1,c)=\Delta(\Run(t,c))$.  This is primitive recursive.  Induction
on $v$ gives the quantifier-free semigroup identity
\[
 \Run(u+v,c)=\Run(v,\Run(u,c)).                                  \tag{3}
\]
Induction also gives absorption of $H$ for any further number of steps.
An injective renaming of the finitely many control labels commutes with
(2); induction on time proves that the dense, native, and symbolic labels
have the same behavior.

\subsection{Finite words and arbitrary finite spare lists}

A finite binary word is coded by its length $\ell$ and its low-bit-first
integer value.  Concatenation satisfies
\[
 |uv|=|u|+|v|,\qquad
 \enc(uv)=\enc(u)+2^{|u|}\enc(v).                                 \tag{4}
\]
The usual recursion defining exponentiation and word append proves (4),
associativity, and reversal of concatenation by quantifier-free induction.
Appending an infinite zero exterior simply uses the same integer
$\enc(u)$, without its finite length.  This is why the interpreter's
exterior-zero absorption is exact equality of configurations, rather than
an unjustified equality of raw untrimmed word lists.

For an explicit coding of arbitrary lists of natural numbers, use a triple
$(\ell,b,N)$ with $b\geq3$, $N<b^\ell$, and all the base-$b$ digits
$d_j$ for $j<\ell$ in $\{1,\ldots,b-1\}$.  The list entries are
$v_j=d_j-1$.  Validity is a bounded primitive recursive test; an arbitrary
finite list has such a code by choosing a sufficiently large base.  Prefix,
suffix, reversal, and the values $v_j$ are primitive recursive.  Appending
the entry 1 replaces $(\ell,b,N)$ by
$(\ell+1,b,N+2b^\ell)$ and preserves validity.  The empty list has
$\ell=N=0$.

Define the block families and their list concatenations by
\[
 B(v)=0\,1^{v+1},\qquad C(v)=1\,0\,1^v,\qquad
 B(X)=\prod_{v\in X}B(v),\quad C(X)=\prod_{v\in X}C(v),
\]
and put
\[
 W(X)=\sum_{v\in X}(v+2)=|B(X)|=|C(X)|.                          \tag{5}
\]
These are primitive recursive functions of the list code.  All inductions
on lists in this note mean quantifier-free induction on the prefix length,
with the full list code, the base, and all outside contexts as free
parameters.  Prefix lengths can be clamped to $\ell$, so the asserted
equation remains unchanged beyond the end of the list.  This requires no
induction on formulas with an unbounded quantifier.

\section{Interpretation of the actual symbolic interpreter}

The coefficient vectors have thirteen slots.  For a valuation
\[
 (1,n,x_0,\ldots,x_8,s,W(X)),                                    \tag{6}
\]
evaluate a vector by its scalar product with (6).  The entries are natural
except that signed coefficients are allowed for the motion vector.
Crucially, slot 12 is \emph{the width of the same list $X$ occurring in the
tape atoms}.  Treating it as an unrelated independent number would be
unsound.  All actual traces use the single identifier
\path{spare_tail}, which is preserved when a list atom moves or changes
family.  The separate one-block derivation reuses slot 2 for a natural $v$.

A run atom $(\texttt{run},b,a)$ denotes $b^{\langle a,\Val\rangle}$.
A list atom denotes $B(X)$ or $C(X)$, with reversal if its flag is set.
The blank atom denotes an infinite zero exterior.  The opaque atoms in a
one-block derivation denote arbitrary outside tape contexts; the
interpreter rejects any attempt to inspect one of them.

A symbolic state contains a control, a left stream $U$, a right stream
$V$, and direction $d\in\{-1,1\}$.  Both streams are listed outward from
a cut.  When $d=1$, the head scans the first cell of $V$; when $d=-1$, it
scans the first cell of $U$.  For a current head coordinate $h$, its
concrete pair of tape halves is
\[
\begin{cases}
(\enc U,\enc V),&d=1,\\
(\lfloor\enc U/2\rfloor,
 (\enc U\bmod2)+2\enc V),&d=-1.
\end{cases}                                                    \tag{7}
\]
If its control is halt, its concrete interpretation is $H$ regardless of
streams.  Let $\Phi$ denote this interpretation.

\paragraph{Well-formed streams.}
The restricted traces in this proof have the following explicit syntactic
invariant.  Each side consists of finitely many run or list atoms followed
by exactly one exterior atom, either blank or one opaque outside context.
An exterior never occurs internally.  Run bits are 0 or 1, their vectors
have thirteen nonnegative integer coefficients, list families are B or C,
and reversal flags are Boolean.  Every list identifier denotes the same
fixed $X$ of (6).  The direction is $+1$ or $-1$.  In a one-block proof
the two exterior contexts may differ and no list atom is needed.
All initial states generated by \path{prove_blocks} and
\texttt{case} have this form.  Normalization only merges or removes
finite atoms or performs (8); it never moves an exterior.  Literal,
unary, and list events remove or transfer finite prefixes and retain
both exteriors.  Expanding a blank inserts a finite 0 before that same
trailing blank.  An attempt to inspect an opaque exterior fails.
Consequently every generated successful trace preserves this invariant.
The soundness proposition below includes it as a hypothesis; it makes
no assertion about malformed arbitrary Python atom lists.

\subsection{Every normalization rule}

Checking a run's bit and nonnegative coefficients ensures that its
valuation is a finite binary run for every valuation (6).  Removing a run
whose entire coefficient vector is zero preserves the word.  A nonzero
vector may still evaluate to zero; retaining that empty run is harmless.
Merging adjacent equal-bit runs uses $b^a b^c=b^{a+c}$.  A finite zero run
beside the infinite zero exterior changes neither binary tape half.
Expanding the blank exterior to $0$ followed by the blank exterior also
changes neither half.  These facts follow directly from (4), or by
induction on the run length.

The only remaining normalization rules use
\[
 C(X)1=1B(X),\qquad 1\rev(C(X))=\rev(B(X))1.                     \tag{8}
\]
Here is the uniform proof, not just a test of short lists.  The empty-list
case is $1=1$.  For one entry,
$C(v)1=10\,1^{v+1}=1B(v)$, using the recursion for a unary run.  For an
appended entry, assuming the identity for a prefix $Y$,
\[
 C(Y)C(v)1=C(Y)1B(v)=1B(Y)B(v).
\]
This is a quantifier-free induction on prefix length using (4).  Reversing
the equality proves the second identity in (8), with reversal proved by
its length recursion.  The interpreter peels one 1 only when the adjoining
run has constant coefficient at least 1, so the remainder is nonnegative
at every valuation.  Each recursive normalization of this kind replaces
one C-list atom by a B-list atom.  Thus at most the initial number of
C-list atoms plus one passes are needed, providing a primitive recursive
bound for normalization itself.

It follows that every call of \texttt{normal} preserves $\Phi$, head
coordinate, and physical time.  This conclusion includes empty lists and
zero-valued symbolic runs.

\subsection{Literal and unary transitions}

For a literal event the interpreter checks a literal row and requires
the run's constant coefficient to be at least 1.  Hence the first bit
exists at every valuation.  Removing it, writing the row's output bit,
and either keeping or reversing direction has exactly the four cases
given by (2) and (7).  Time increases by 1 and head coordinate by $d$.
The proof is substitution in those equations; it requires no induction.

For a unary event the checked row has the same entry and exit control
and moves in the current direction.  For either bit written by that row,
crossing a run of length $a$ in that direction takes exactly $a$ literal
steps, writes a run of the output bit, and transfers it to the opposite
stream.  The induction assertion after $j\leq a$ steps is that the
unprocessed source suffix has length $a-j$, the written destination
prefix has length $j$, the control is unchanged, and the head displacement
is $dj$.  For the next step this assertion follows from (2); for $a=0$
it is the identity.  All quantities and the assertion are primitive
recursive and quantifier free.  Substituting an affine run length for
$a$ proves exactly the interpreter's unary rule.

\section{The four block and list traversal lemmas}

The following are the actual literal rows used by the four one-block
derivations.  Labels are the interpreter's symbolic labels, not an
additional machine.

\begin{center}
\begin{tabular}{llll}
\toprule
Entry / direction / action & First row & Second row & Loop row\\
\midrule
73 / R / B to B & $73,0\mapsto0R83$ & $83,1\mapsto1R73$
 & $73,1\mapsto1R73$\\
95 / R / B to C & $95,0\mapsto1R109$ & $109,1\mapsto0R95$
 & $95,1\mapsto1R95$\\
48 / L / B to B & $48,1\mapsto1L63$ & $63,0\mapsto0L48$
 & $63,1\mapsto1L63$\\
112 / L / C to B & $112,0\mapsto1L131$ & $131,1\mapsto0L112$
 & $112,1\mapsto1L112$\\
\bottomrule
\end{tabular}
\end{center}

For entry 73 from the left, two literal steps cross the initial 0 and
first 1 of $B(v)$; the loop crosses the remaining $v$ ones.  The block is
unchanged.  For entry 95, the first two steps write 1 and 0 instead, and
the $v$ remaining ones are unchanged, giving $C(v)$.  For entry 48 from
the right, one step crosses the last 1, the loop crosses the other $v$
ones, and the final step crosses the 0; the result is $B(v)$.  For entry
112 from the right of $C(v)$, the loop crosses the $v$ terminal ones;
two steps then replace the 0 by 1 and the initial 1 by 0, giving $B(v)$.
The unary lemma proves the loop in each description for every natural
$v$, including zero.

In all four cases the exact duration is $v+2$, the exit control equals
the entry control, and the direction is unchanged.  The displacement is
$d(v+2)$.  No cell outside the block is inspected or written: the last
step merely moves onto the next outside cell.  These statements hold for
arbitrary left and right contexts, because every read has been accounted
for within the displayed block.  They are the four generic derivations
in \path{list_lemmas.json}; \path{prove_blocks} rederives them
from the current literal table rather than trusting state names.

\begin{lemma}[Uniform list traversal]
For each of these four entries, replace the single block by a finite list
of blocks, read in the specified direction.  The machine crosses the
whole list, replaces each block by the displayed output family, returns
to the same control and direction, takes exactly $W(X)$ steps, and has
displacement $dW(X)$.  It leaves both outside contexts unchanged.
\end{lemma}

\begin{proof}
For rightward traversal, after $j$ blocks the processed prefix is the
output-family word for the first $j$ entries, the unprocessed suffix is
the input-family word for the remaining entries, the control is the
entry control, and the elapsed time is
$\sum_{i<j}(v_i+2)$.  This is an equality of encoded complete
configurations with the untouched contexts as free parameters.  The
one-block calculation and (3) establish the next case.  The case $j=0$
is the identity.  Quantifier-free induction on $j$, clamped to the list
length as above, proves the claim.  For leftward traversal use the same
invariant with entries in reverse order.  The total is still $W(X)$.
Each step preserves both outside contexts, so the final equality does
also.  No bound on the list entries or the number of entries is used.
\end{proof}

This lemma proves exactly the symbolic list event: transfer the opaque
list atom to the other side of the cut, toggle its reversal flag,
replace its family as specified, add slot 12 to the clock, and add
$d$ times slot 12 to the motion.  It also proves the zero-duration
empty-list case.  A whole source instruction has positive duration
independently of this possible zero-duration event.

\section{Certificate soundness inside \texorpdfstring{$\ipra$}{iPRA}}

The restricted certificate interpreter has a primitive recursive coding.
Atoms, affine coefficient vectors, and finite event lists are finite
codes.  Literal row lookup is finite.  Coefficient arithmetic, sign tests,
and comparisons are primitive recursive.  Normalization has the bound
proved above.  The one-block derivation is bounded by 50 symbolic events;
a source case is bounded by 4000 events.  A failed check returns a
distinguished failure code.  A successful event checks precisely one
of the rules just proved.  The sole list identifier is bound to the same
list code throughout, and its width is defined by (5), not supplied as
an unconstrained oracle.

\begin{proposition}[Soundness of the implemented event calculus]
For any successful restricted trace with well-formed symbolic start $A$
as above, end $B$, clock vector $a$, and motion vector $m$, and any
valuation (6) with a valid finite list code, the literal machine satisfies
\[
 \Run(\langle a,\Val\rangle,\Phi(A,h))
   =\Phi(B,h+\langle m,\Val\rangle).                             \tag{9}
\]
This proposition is provable in $\ipra$.
\end{proposition}

\begin{proof}
Use quantifier-free induction on the trace prefix length.  The induction
formula is the implication from successful-prefix checking, the initial
well-formedness test, and valid parameters to well-formedness and the
configuration equality (9) for that prefix.  All its tests and
configurations are primitive recursive.  The empty prefix
uses $\Run(0,c)=c$.  Normalization preserves interpretation by (8) and
the preceding elementary rules.  A literal event uses (2), a unary event
uses the unary induction, and a list event uses the proved list lemma.
Equation (3) composes the new segment with the prior prefix, while
coefficient addition gives exactly the new clock and displacement.
The syntactic preservation argument above proves the well-formedness
part of the induction step.
Thus the induction step proves (9) for every type of event accepted by
the interpreter.  There is no additional symbolic execution rule.
\end{proof}

For these machines the successful restricted checks are concrete finite
computations.  Their closed arithmetic instances are provable by
unfolding the defining primitive recursions on numerals.  Applying the
proposition gives their universal instances; the parameters remain free
throughout the derivation.  This is the precise division of labor:
the finite computations decide which literal derivations were accepted,
and the arithmetic proofs above make each accepted derivation valid for
arbitrary natural registers, gap, and finite list.  The stored
\texttt{contracts.json} contains the durations and event counts, not
the entire event traces; the bounded literal replay regenerates the
traces, and its \path{keep_trace} option can expose them.  A receipt
alone, without that replay and the rule proofs, is not a universal
simulation theorem.

No proof-assistant kernel has checked this derivation.  The claim here
is an ordinary written proof of an $\ipra$ theorem using the accompanying
finite computations.  The uniform arguments are actual quantifier-free
induction arguments; they do not assume that a computer-tested property
extends to arbitrary inputs.

\section{All source frames and their closure}

Let the logical registers be $R_0,\ldots,R_8$.  The physical permutation
is
\[
 \pi=(6,2,0,4,3,7,8,5,1),\qquad V_{\pi(i)}=R_i.                 \tag{10}
\]
Let $p$ be the binary address of a source site, let $n\geq0$, and let
$Q=[s]X$ be an arbitrary nonempty finite list.  At the boundary control
84 (96 in the symbolic names), with head at $h=-2$, the canonical frame
has outward left stream $1$ followed by blank zeros and right stream
\[
 p\;0^{48+n-|p|}\;1\;
 B(V_0)\cdots B(V_8)\; B(s)B(X)\;0^\infty.                     \tag{11}
\]
This is an exact complete tape, including both infinite exteriors.
The maximum address length is 37 or 36, so every initial zero-gap
exponent in (11) is nonnegative for every $n$.

For an increment the target register is increased by 1 and the spare
list becomes $Q[1]$.  For a positive decrement the target register is
decreased by 1 and $n$ becomes $n+1$.  A zero decrement, break, or
continue leaves registers, $n$, and $Q$ unchanged.  Halting yields $H$.
The nonempty-list invariant is preserved, and the explicit base-$b$
encoding supports $Q[1]$ without recoding the other entries.

The source case splits a decrement into its zero case and its positive
case $R_i=x_i+1$.  Every natural input is covered: decidability of zero
and the identity $R_i>0\Rightarrow R_i=\operatorname{pred}(R_i)+1$
are quantifier-free arithmetic facts.  Other registers are independent
arbitrary naturals.  In a positive case the corresponding valuation
slot is $\operatorname{pred}(R_i)$; the other logical slots have their
register values.  In the zero case that register is literally zero in
the symbolic frame.

The interpreter computes a binary successor prefix using increment
and cut operations.  A priori this prefix need not itself be a source
leaf.  The necessary endpoint assertion is equality after padding with
zeros, not an unproved identification of prefixes.  For every one of
the 4777 nonhalting outcomes in each machine, the additional finite audit
checks
\[
 p_{\rm computed}\,0^{48-|p_{\rm computed}|}
  =p_{\rm actual\ successor}\,0^{48-|p_{\rm actual\ successor}|},
 \qquad \mathrm{overflow}=0.                                    \tag{12}
\]
The actual successor comes from the graph that agrees with
\texttt{instructions.json}; the arithmetic and interval resolvers agree
on all edges.  The full prefix-tree syntax is checked too.  Appending
the same extra $n$ or $n+1$ zeros preserves (12), so it proves closure
of the universal target frame (11).  The literal certificates give zero
motion in all these 4777 cases.  Thus the boundary head remains at
$h=-2$; there is no unaccounted head translation.

Each duration vector has nonnegative coefficients, with constant at
least 3.  The componentwise maximum of the vectors in either machine is
\[
 (534,8,4,4,4,4,4,4,2,4,4,2,2).                                \tag{13}
\]
Consequently the exact duration $D$ selected by a source state is
primitive recursive and, before halt, satisfies
\[
 3\leq D\leq534+8n+4\sum_{i=0}^{8}R_i+2W(Q).                  \tag{14}
\]
For a positive decrement its slot value is at most $R_i$; for the spare
slots, $2s+2W(X)\leq2W(Q)$.  These observations justify (14) directly
from (13).  The exact duration, not this upper bound, is used in the
simulation.  The sole halt contract has constant duration 3 and reaches
native control 254.  Its halted tape is not required by (11).

\begin{proposition}[Universal local contracts]
Let $F$ be the extended source step just described, including the gap
and spare-list bookkeeping, and let $E$ map a nonhalting source state
to (11), with $E(H)=H$.  Put $D(H)=0$.  For every valid extended source
state $z$, $\ipra$ proves
\[
 \Run(D(z),E(z))=E(F(z)).                                       \tag{15}
\]
Moreover $F$ preserves validity, fixes $H$, and $D(z)\geq3$ for
$z\neq H$.
\end{proposition}

\begin{proof}
For each valid program counter choose its fixed contract, splitting
decrement by its decidable zero test.  Apply (9), (10)--(12), and the
observed coefficient facts (13).  These are finitely many cases, covering
all 3671 program counters and 4778 outcomes.  The halt case uses the
canonical physical halt code, so it also has exact endpoint equality.
For $H$, (15) is $\Run(0,H)=H$.  The register changes preserve natural
values, the gap stays nonnegative, and append preserves nonempty valid
spare lists.  This proves every assertion.
\end{proof}

\section{Startup and composition with primitive recursive clocks}

The independent startup simulation uses the literal native rows, a
20,000-cell zero array, and origin 10,000.  It checks at every one of
57922 transitions that the head is in the array and the control is
nonhalting.  At the endpoint the control is 84, the head coordinate is
$-2$, and the \emph{entire} array equals (11) for site 25, $n=0$,
logical registers $(2,0,\ldots,0)$, and 26 spare entries all equal to 1.
Outside the array no cell was visited and the tape remains zero.
Hence this finite computation establishes the whole-tape equality
\[
 \Run(T_0,c_{\rm blank})=E(z_0),\qquad T_0=57922.                \tag{16}
\]
This closed computation is an arithmetic numeral fact.  The nonhalting
endpoint and absorbing-halt convention also imply that no earlier
startup configuration was halted.

Define, by primitive recursion,
\[
 z_k=F^k(z_0),\qquad
 \tau(0)=0,\qquad \tau(k+1)=\tau(k)+D(z_k).                      \tag{17}
\]
The source machine has a fixed finite instruction array, primitive
recursive register updates, and the explicit list operations above.
Thus $z_k$ and $\tau(k)$ are primitive recursive functions of $k$.
There is no minimization or appeal to termination of an unbounded search.

Equations (3), (15), and (16) give, by quantifier-free induction on $k$,
\[
 \Run(T_0+\tau(k),c_{\rm blank})=E(z_k).                         \tag{18}
\]
The induction formula includes the characteristic validity test for
$z_k$, which is preserved by (15).  If $z_k=H$, then $z_{k+1}=H$ and
$D(z_k)=0$; (18) is still exact.  This explicitly handles all times
after source halt.

A second quantifier-free induction proves
\[
 z_k\neq H\ \Longrightarrow\ \tau(k)\geq3k.                    \tag{19}
\]
Indeed, if $z_{k+1}\neq H$, absorption of the source halt code implies
$z_k\neq H$, and (14) increases the clock by at least 3.  The base case
is immediate.

If the source is halted at $k$, (18) gives physical halt at the explicit
primitive recursive time $T_0+\tau(k)$.  Conversely suppose physical
halt has occurred at time $t$, and set $k=t+1$.  If $z_k\neq H$, then
(19) gives $T_0+\tau(k)>t$.  Absorption of physical halt would make the
left side of (18) equal to $H$, whereas its right side is a nonhalting
boundary frame at control 84, a contradiction.  Equality with the fixed
halt code is decidable in $\ipra$, so this proves $z_{t+1}=H$
intuitionistically.  It supplies the explicit reverse witness $t+1$;
no Markov principle, excluded middle for arbitrary formulas, or
unbounded least-time principle is used.

Finally, erasing $n$ and $Q$ from the extended state gives exactly the
ordinary program state $(\text{site},R_0,\ldots,R_8)$ in
\texttt{instructions.json}.  The erasure commutes with one source step,
including its canonical halt code.  Quantifier-free induction on $k$
therefore equates the projected extended run with the ordinary register
run.  This proves (1).

\section{Scope of the physical audit and its composition}

The actual interpreter has no additional unproved all-input abstraction:
normalization is finite-word equality; unary execution is length
induction; the four opaque-list events are the exact block calculations
and prefix inductions above; literal events are literal transitions.
Every segment duration is an explicit primitive recursive affine
expression with the stated list-width interpretation.  The complete
source-frame equality, positive-duration facts, exact startup, and
composition close the physical/source bridge.

Because the proofs are $\ipra$ derivations with arbitrary parameters,
their universal conclusions apply within every model of $\ipra$.
They are not limited to standard finite values sampled by Python.
The code evaluates the fixed symbolic certificates; induction supplies
their interpretation at arbitrary values, including model-internal
finite lists.

The physical certificate alone does not identify its register program
with a PA proof search. The companion sections in the complete report
supply the exact source-to-command implementation, the command compiler,
the current Foundation PA adapter, and the numerical soundness and
completeness constructions. They compose the requested
$\operatorname{Con}(\mathsf{PA})$ equivalence in iPRA using the stated
weak-arithmetic transfer theorems. The present note proves the physical
interface used by that composition. It is a written arithmetic proof,
not a claim that a new end-to-end Lean theorem has been built.

\paragraph{Reproduction.}
Run \path{verification/check_literal.py} on each machine directory for
full literal replay. Run \path{verification/physical_ipra_audit.py}
with \texttt{--bundle PATH --output FILE} for this note's additional
checks. Keep Python assertions enabled. The delivered audit receipt,
including input hashes, is
\path{verification/physical_ipra_audit_receipt.json}.

\end{document}
